\section{Construction and Implementation}
\label{app:construction}

This appendix records the predicate features, bank construction and
checking choices, and training details used by the method.

\begin{table}[t]
\centering\small
\caption{Frozen banks used in the reported experiments. Each originates
from 36 generation consultations; checking requests are additional.
The two MultiRoom banks share the same initial consultations.}
\label{tab:banks}
\begin{tabular}{@{}llllr@{}}
\toprule
Task & Student & Check & Action filter & Rules\\
\midrule
DoorKey & Both & Blind & None & 16\\
MultiRoom & PPO & None & None & 32\\
 & Count-PPO & Blind & None & 7\\
KeyCorridor & Both & Self-check & Before freezing & 28\\
Crafter & Both & Self-check & Runtime & 24\\
\bottomrule
\end{tabular}
\end{table}

\subsection{Predicate Features and Progress}
\label{app:predicate_features}

The final vocabularies are frozen before evaluated training. MiniGrid observation predicates describe the front
cell, carried object, visibility of task-relevant object types, coarse
direction and offsets to their closest visible instances, and visible
door states. The carried object is encoded at the agent's cell in the
symbolic observation. An unseen object's direction and offsets are
$\mathsf{unknown}$.
Crafter predicates describe the facing cell and direction, daylight,
nearby table and furnace, distance and first move towards the closest
visible instance of each of 13 object types, and binned inventory,
tools and vital statistics.
Rule conditions use values from these finite domains. An omitted
condition imposes no constraint; $\mathsf{unknown}$ is a runtime
information state, not a value requested by a generated rule.

KeyCorridor's executor computes \texttt{door\_unlocked} by scanning
the simulator grid for remaining locked doors. In the evaluated task,
the student is the only agent that can unlock the door, and unlocking
cannot be reversed. A separate software check reconstructed this flag
from executed actions and local observations and agreed with the
grid-derived flag at every compared state. That reconstruction was
not used by the training implementation and does not establish what
the learned policy remembers. Crafter's executor reads its 22 progress
facts from the game's episodic achievement counters. Neither set of
progress facts is added to the policy input.

\subsection{Generation and Checking}
\label{app:bank_construction_details}

\paragraph{Generation states.}
MiniGrid uses 36 distinct construction layouts and one example state
per layout. DoorKey allocates 12 examples to each phase: key on the
floor, key in hand, and door open, with the agent placed at a randomly
chosen reachable cell. MultiRoom and KeyCorridor allocate 12 examples
to each of the early, middle and late thirds of a planner trajectory.
A short random walk is added with probability $1/2$ after reaching
the sampled point.
The planner supplies states, not training action targets.
Crafter's initial examples use short random walks and constructed
changes to inventories and crafting stations. Initial calls each
contain a single prompt and no previous model response. We use
\texttt{gpt-5-mini-2025-08-07}, low reasoning effort and structured output.
The rule proposals from these 36 calls are collected into one initial
bank. The count of 36 is a fixed generation budget, not an established
minimum or optimum. Development compares whole-bank variants; it
does not train one student for each individual proposed rule.

\paragraph{Checking states.}
MiniGrid checking pools contain additional states from the layouts
used for generation, so they need not be the original consultation
states. A checking request contains up to five states where the rule
applies. Blind checks ask for an action at each state without showing
the rule. Self-checks show the rule and can revise its conditions,
exceptions or action, or withdraw it. Their treatment of rules with
no applicable states is specified in the main method.
Crafter's later checking pools also contain states visited by
unadvised development policies and random play. Its selected bank
rechecks earlier self-checked rules using the extended progress
vocabulary. Actions with the same measured effect as a no-op are
flagged for the LLM using separate simulator copies.

\paragraph{Pooling and action checks.}
In KeyCorridor, checks for identical proposed rules are pooled. An
original rule remains unchanged only if every copy's check leaves it
unchanged; otherwise, distinct refinements are retained and withdrawn
copies contribute no rule.
Fixed action checks are then compiled into rule conditions. For
example, \texttt{pickup} permits a key or ball in front with empty
hands, whereas \texttt{drop} permits a carried key with empty floor
in front. A missing condition is added by expanding into allowed
cases; incompatible cases are discarded. These are specified action
restrictions, not a complete guarantee that every retained action
changes the environment. Crafter applies its filter at runtime using
its action requirements and recipe table, excluding \texttt{noop}.

\paragraph{Development selection.}
DoorKey and MultiRoom compare checked and unchecked candidate banks
using mean teacher-free learning-curve area on development runs,
separately for PPO and Count-PPO. The selection procedure chooses
no advice if neither candidate improves on its development control.
The chosen MultiRoom variants are the unchecked bank for PPO and
the checked subset for Count-PPO. These choices precede the reported
evaluation runs. KeyCorridor and Crafter additionally use direct bank
execution with random actions when the executor abstains; Crafter
includes development learning runs. MiniGrid optimal-action planners
provide held-out development label-precision diagnostics for choosing
checking variants. Their actions never enter a frozen bank or the
student's training targets.

\subsection{Learning Details}
\label{app:learning_details}

\paragraph{PPO objective.}
The ordinary PPO loss~\citep{schulman2017ppo} is
\begin{align}
\mathcal{L}_{\mathrm{PPO}}={}&-\hat{\mathbb{E}}_{\mathcal{M}}
\bigl[\min\bigl(\rho_t\hat{A}_t,
\operatorname{clip}(\rho_t,1-\epsilon,1+\epsilon)\hat{A}_t\bigr)\bigr]
\nonumber\\
&+c_V\mathcal{L}_V-c_H\hat{\mathbb{E}}_{\mathcal{M}}
\bigl[\mathcal{H}(\pi_\theta(\cdot\mid x_t))\bigr],
\label{eq:ppo}
\end{align}
where $\rho_t=\pi_\theta(a_t\mid x_t)/
\pi_{\theta_{\mathrm{old}}}(a_t\mid x_t)$,
$\hat{A}_t$ is an advantage estimate, $\mathcal{L}_V$ the value loss,
and $\mathcal{H}$ policy entropy. Every executed action is sampled
by the student, so this remains the ordinary on-policy ratio.

\paragraph{Count-PPO.}
Counts reset each episode. In MiniGrid, they index the local
$7\times7$ observation. The intrinsic bonus is normalized using
a running scale and has a separate value head; its advantage is
combined with the task advantage as
$\hat{A}_t=\hat{A}^{\mathrm{task}}_t+
\beta\hat{A}^{\mathrm{int}}_t$.
Both value losses enter $\mathcal{L}_V$.
In Crafter, counts index the $9\times7$ grid observation excluding
inventory, and the bonus is added to the reward of a single value
head. Evaluation reports task metrics without a count bonus.

\paragraph{Schedule and algorithm.}
$\lambda(\cdot)$ denotes the schedule function;
$\lambda_k=\lambda(\tau_k)$ is its value for rollout iteration $k$,
where $\tau_k$ is the completed training fraction at that iteration's start.
The initial weight and floor, $\lambda_0$ and $\lambda_{\min}$, are
fixed hyperparameters. Algorithm~\ref{alg:training} uses $\lambda_k$
for both target generation and the minibatch objective.
When advice is disabled, $y_t=\bot$ and $m_t=0$; only transitions
with $m_t=1$ enter the advice loss.

% Requires algorithm and algpseudocode. Included by the Appendix file.
\begin{algorithm}[t]
\caption{Training with executable advice. Count-PPO additionally stores
the count bonus; MiniGrid uses a separate intrinsic reward stream.}
\label{alg:training}
\begin{algorithmic}[1]
\Require frozen bank $\mathcal{B}$, schedule $\lambda(\cdot)$, $K$ iterations
\For{$k=1,\dots,K$}
  \State $\tau_k\gets$ completed training fraction
  \State $\lambda_k\gets\lambda(\tau_k)$
  \For{each rollout step $t$ in every environment}
    \State sample $a_t\sim\pi_\theta(\cdot\mid x_t)$
    \State $y_t\gets E_{\mathcal{B}}(z_t)$ if $\lambda_k>0$; otherwise $\bot$
    \State $m_t\gets\mathbb{I}[y_t\ne\bot]$
    \State execute $a_t$; store transition, $y_t$ and $m_t$
  \EndFor
  \State compute advantages $\hat{A}_t$
  \For{each minibatch $\mathcal{M}$ and epoch}
    \State update $\theta$ using
    $\mathcal{L}_{\mathrm{PPO}}+\lambda_k\mathcal{L}_{\mathrm{adv}}$
  \EndFor
\EndFor
\State \Return $\pi_\theta$ \Comment{acts alone at evaluation}
\end{algorithmic}
\end{algorithm}

